\documentclass[10pt,a4paper]{amsart}
\usepackage{amsmath,amssymb,amsthm,graphicx}
% Retained minimal preamble for the Stacks Project chapter sites-cohomology.tex (Cohomology on Sites).
% Verbatim excerpt of the chapter preamble: Xy-pic setup (native preamble lines 14-18) and the native
% theorem declarations with the mathematical macros (native preamble lines 154-217). Copied mechanically;
% nothing added, removed or redefined. The wrapper class and the amsmath/amssymb/amsthm/graphicx base
% packages are supplied by the assembler, so the class/externaldocument part of the native preamble is omitted.
% For commutative diagrams we use Xy-pic
\usepackage[all]{xy}

% We use 2cell for 2-commutative diagrams.
\xyoption{2cell}
% Theorem environments.
%
\theoremstyle{plain}
\newtheorem{theorem}[subsection]{Theorem}
\newtheorem{proposition}[subsection]{Proposition}
\newtheorem{lemma}[subsection]{Lemma}

\theoremstyle{definition}
\newtheorem{definition}[subsection]{Definition}
\newtheorem{example}[subsection]{Example}
\newtheorem{exercise}[subsection]{Exercise}
\newtheorem{situation}[subsection]{Situation}

\theoremstyle{remark}
\newtheorem{remark}[subsection]{Remark}
\newtheorem{remarks}[subsection]{Remarks}

\numberwithin{equation}{subsection}

% Macros
%
\def\lim{\mathop{\mathrm{lim}}\nolimits}
\def\colim{\mathop{\mathrm{colim}}\nolimits}
\def\Spec{\mathop{\mathrm{Spec}}}
\def\Hom{\mathop{\mathrm{Hom}}\nolimits}
\def\Ext{\mathop{\mathrm{Ext}}\nolimits}
\def\SheafHom{\mathop{\mathcal{H}\!\mathit{om}}\nolimits}
\def\SheafExt{\mathop{\mathcal{E}\!\mathit{xt}}\nolimits}
\def\Sch{\mathit{Sch}}
\def\Mor{\mathop{\mathrm{Mor}}\nolimits}
\def\Ob{\mathop{\mathrm{Ob}}\nolimits}
\def\Sh{\mathop{\mathit{Sh}}\nolimits}
\def\NL{\mathop{N\!L}\nolimits}
\def\CH{\mathop{\mathrm{CH}}\nolimits}
\def\proetale{{pro\text{-}\acute{e}tale}}
\def\etale{{\acute{e}tale}}
\def\QCoh{\mathit{QCoh}}
\def\Ker{\mathop{\mathrm{Ker}}}
\def\Im{\mathop{\mathrm{Im}}}
\def\Coker{\mathop{\mathrm{Coker}}}
\def\Coim{\mathop{\mathrm{Coim}}}

% Boxtimes
%
\DeclareMathSymbol{\boxtimes}{\mathbin}{AMSa}{"02}

%
% Macros for moduli stacks/spaces
%
\def\QCohstack{\mathcal{QC}\!\mathit{oh}}
\def\Cohstack{\mathcal{C}\!\mathit{oh}}
\def\Spacesstack{\mathcal{S}\!\mathit{paces}}
\def\Quotfunctor{\mathrm{Quot}}
\def\Hilbfunctor{\mathrm{Hilb}}
\def\Curvesstack{\mathcal{C}\!\mathit{urves}}
\def\Polarizedstack{\mathcal{P}\!\mathit{olarized}}
\def\Complexesstack{\mathcal{C}\!\mathit{omplexes}}
% \Pic is the operator that assigns to X its picard group, usage \Pic(X)
% \Picardstack_{X/B} denotes the Picard stack of X over B
% \Picardfunctor_{X/B} denotes the Picard functor of X over B
\def\Pic{\mathop{\mathrm{Pic}}\nolimits}
\def\Picardstack{\mathcal{P}\!\mathit{ic}}
\def\Picardfunctor{\mathrm{Pic}}
\def\Deformationcategory{\mathcal{D}\!\mathit{ef}}
\usepackage{cleveref}
\crefname{originalreference}{original reference}{original references}
\Crefname{originalreference}{Original reference}{Original references}
\makeatletter
\newlabel{lemma-existence-derived-lower-shriek}{{07AC}{source}{}{sourcepointer-1}{}}
\newlabel{lemma-existence-derived-lower-shriek@cref}{{[originalreference][0][]07AC}{[1][0][]source}{}{}{}}
\makeatother
\crefname{originalreference}{original reference}{original references}
\Crefname{originalreference}{Original reference}{Original references}
\makeatletter
\newlabel{lemma-pullback-injective-pre-limp}{{0D6X}{source}{}{sourcepointer-2}{}}
\newlabel{lemma-pullback-injective-pre-limp@cref}{{[originalreference][0][]0D6X}{[1][0][]source}{}{}{}}
\makeatother
\title{Cohomology on an object agrees with cohomology on its image under a continuous cocontinuous functor}
\author{The Stacks Project Authors}
\date{stacks-project:sites-cohomology.tex:local-snapshot-20261011; GNU FDL 1.2 or later}
\begin{document}
\maketitle
\noindent\textit{Editorial note.} All delivered mathematical text is the author's own, copied from the pinned Stacks Project chapter file sites-cohomology.tex; no statement or proof step was rewritten, reordered or removed.\par\medskip
\noindent\textit{Editorial note.} The counted proof is the author's complete argument: the proof is written out in full: the bounded below case with Leray's acyclicity lemma, the reduction of (1) to (2), and the displayed adjunction chain for the unbounded case.\par\medskip
\noindent\textit{Editorial note.} The label lemma-pullback-same-cohomology names the counted statement, and the proof environment immediately following it in the source is delivered verbatim as its proof.\par\medskip
\noindent\textit{Editorial note.} The counted statement is carried by the author's own lemma environment in the chapter source.\par\medskip
\noindent\textit{Editorial note.} The statement and proof are copied from the pinned chapter file; the two cited inputs are kept as external pointers with their canonical Stacks tags.\par\medskip
\noindent\textit{Allowed standard prerequisites (editorial).}\begin{itemize}
\item Acyclicity of the pullback of an injective module. Same chapter, the lemma stating that the pullback along a continuous cocontinuous functor of an injective module has no higher cohomology on the objects of the source site. Supplies the input for Leray's acyclicity lemma in the bounded below case.
\item Existence of the derived lower shriek and its adjunction. Same chapter, the lemma constructing the derived lower shriek as a left adjoint to the pullback. Supplies the middle equality of the displayed chain in the unbounded case.
\end{itemize}\medskip
\ifdefined\hypertarget\hypertarget{sourcepointer-1}{}\fi
\noindent\textit{Source reference.} Original reference 07AC: sources/stacks-sites-cohomology/pinned-original-source.tex, source lines 9310--9310. Prerequisite pointer invoked at proof line(s) 9539; it names a standard fact of the same chapter that the counted argument uses, and every step of the counted proof built on it is retained verbatim in the delivered body.\par\medskip
\ifdefined\hypertarget\hypertarget{sourcepointer-2}{}\fi
\noindent\textit{Source reference.} Original reference 0D6X: sources/stacks-sites-cohomology/pinned-original-source.tex, source lines 9278--9278. Prerequisite pointer invoked at proof line(s) 9528; it names a standard fact of the same chapter that the counted argument uses, and every step of the counted proof built on it is retained verbatim in the delivered body.\par\medskip
\begin{lemma}
\label{lemma-pullback-same-cohomology}
Let $u : \mathcal{C} \to \mathcal{D}$ be a continuous and cocontinuous
functor of sites. Let $g : \Sh(\mathcal{C}) \to \Sh(\mathcal{D})$
be the corresponding morphism of topoi. Let $U \in \Ob(\mathcal{C})$.
\begin{enumerate}
\item For $M$ in $D(\mathcal{D})$ we have
$R\Gamma(U, g^{-1}M) = R\Gamma(u(U), M)$.
\item If $\mathcal{O}_\mathcal{D}$ is a sheaf of rings and
$\mathcal{O}_\mathcal{C} = g^{-1}\mathcal{O}_\mathcal{D}$, then
for $M$ in $D(\mathcal{O}_\mathcal{D})$ we have
$R\Gamma(U, g^*M) = R\Gamma(u(U), M)$.
\end{enumerate}
\end{lemma}

\begin{proof}
In the bounded below case (1) and (2) can be seen by representing
$K$ by a bounded below complex of injectives and using
Lemma \ref{lemma-pullback-injective-pre-limp} as well as
Leray's acyclicity lemma.
In the unbounded case, first note that
(1) is a special case of (2). For (2) we can use
$$
R\Gamma(U, g^*M) =
R\Hom_{\mathcal{O}_\mathcal{C}}(j_{U!}\mathcal{O}_U, g^*M) =
R\Hom_{\mathcal{O}_\mathcal{D}}(j_{u(U)!}\mathcal{O}_{u(U)}, M) =
R\Gamma(u(U), M)
$$
where the middle equality is a consequence of
Lemma \ref{lemma-existence-derived-lower-shriek}.
\end{proof}

\end{document}
